\section{Conclusion}
\label{sec:conclusion}

Temporary industrial support is often evaluated by asking whether the supported activity survives once the policy is withdrawn. This paper has argued that the question should be split into stages.

First, support can solve an activation problem without solving a capability problem. Second, operating under support can accumulate durable capability, but finite graduation occurs only if the supported operating regime can generate a long-run capability level above the unsupported viability threshold. In the baseline model, this condition is $\lambda/\delta>q^*$, and the implied minimum support duration follows in closed form. Third, surrounding productive structure matters because it can raise capability conversion and lower unsupported thresholds. In a fixed active network, these interactions create a coupled capability system whose finite steady-state condition is $\rho(M)<1$.

Fourth, graduation must be distinguished from two different kinds of dependence. The unsupported dependency closure asks whether the productive configuration survives the instant extraordinary support is removed after support-dependent neighbors are recursively eliminated. Dynamic unsupported viability asks a harder question: whether the surviving configuration reproduces its capability requirements through time. Under a fixed support-free configuration, a sufficient condition is the forward-invariance requirement
\[
M_0\vartheta+b_0\ge\vartheta.
\]
At the threshold of viability, the unsupported productive configuration must reproduce at least its own capability requirements.

Fifth, technological graduability does not imply that firms will choose graduation as their preferred survival route. The playing-field extension distinguishes credibility that the announced learning window will be honored from credibility that the announced terminal condition cannot be privately renegotiated. When the optimized value of graduation $G(\sigma)$ exceeds the optimized value of political dependence $D(\chi)$, graduation is incentive-compatible in the reduced-form sense developed here. The point is not that credible withdrawal mechanically forces investment. It is that political preservation of support can otherwise remain a substitute for becoming capable enough to leave support.

The distinction between durable capability, current dependence, and causal provenance remains central. A supported neighbor may generate learning that remains after the neighbor disappears; that durable learning legitimately enters $q$. A complement that must remain present belongs instead in the current environment. And the question of how much capability support caused belongs to a separate counterfactual attribution exercise.

Finally, graduation is not the terminal developmental criterion. Mature productive relations can and should sometimes disappear. The relevant question is what happens to the capabilities and relations around them. Exit can represent failed graduation, destructive de-accretion, competitive displacement, or successful supersession. The distinction matters because industrial survival and industrial death are both ambiguous signals.

The institutional and structural claims can therefore be summarized together:
\begin{quote}
\textbf{Time under protection is not learning.}
\end{quote}
\begin{quote}
\textbf{Stable rules need not mean constant policies; they mean credible rules governing change.}
\end{quote}
\begin{quote}
\textbf{Industrial survival is not developmental success, and industrial death is not capability death.}
\end{quote}
The central developmental claim remains:
\begin{quote}
\textbf{Development succeeds when productive capability can outlive the productive configurations that created it.}
\end{quote}
These claims are proposed as a bounded research framework, not as established empirical laws or a general theory of development. Their value depends on whether capability stocks, political-dependence routes, dependency closures, dynamic viability conditions, transfer channels, provenance counterfactuals, and de-accretion cascades can be identified in actual historical and contemporary episodes.
